%!TEX root = ../student_thesis.tex

\section*{Abstract}
% Kurzzusammenfassung / Abstract (max. eine Seite)
Electric machines operate at relatively low frequencies yet are typically driven by high-frequency power-electronic circuitry, so they are most efficiently time-integrated at different rates. Their differing character – one a complex machine, modelled as a distributed device with differential equations and simulated numerically using methods such as the Finite Element Method (FEM), the other a circuit solved by its own solver – motivates a partitioned co-simulation rather than a monolithic solve. This work couples the domains with a Waveform Relaxation (WR) scheme, in which an inductance-based reduced-order model of the field serves as the interface condition and, as a linearisation term, accelerates convergence. The scheme is implemented with Xyce, Sandia's open-source SPICE simulator, leveraging its checkpoint/restart feature to re-solve each time window across WR iterations. The implementation is verified against monolithic solutions for cases with varying circuit complexity and a dummy FEM solver. A parameter study over [window length/linearisation update / …] shows that [key result], from which [a concrete guideline, e.g. for choosing the window size] is derived.
